\usepackage{ifthen}
% ======================================================
%   COLORS
%colors
\definecolor{CScolor}{rgb}{1.0, 0.0, 0.0} % Chern Simons
\definecolor{BFcolor} {RGB} {0, 158, 96} % BF
\definecolor{FIcolor}{RGB}{250, 95, 85} % FI
%\colorlet{BFcolor}{BFcolor} % BF

\definecolor{bcolor}{RGB}{4, 55, 242} % bosons
\definecolor{fcolor}{RGB}{255, 36, 0} % fermions
\definecolor{labelcolor}{HTML}{808080} % labels



% ======================================================
%   SYMBOLS SHORTCUTS
\def\elephi{\varphi}
\def\elepsi{\psi}
\def\magphi{\phi}
\def\magpsi{\psi}

\def\with{\text{ with }}
\def\mH{m_H}
\def\Nt{\widetilde{N}}

\def\dir{\mathcal{D}}
\def\neu{\mathcal{N}}

\def\mon{\mathfrak{M}}

\newcommand{\zz}{\mathbb{Z}}
\newcommand{\rr}{\mathbb{R}}
\newcommand{\cc}{\mathbb{C}}
\newcommand{\nn}{n}
\newcommand{\NN}{\mathcal{N}}
\def\dualto{
	\qquad\leftrightarrow\qquad
}
\def\CSg{
	\text{CS$_g$}
}

\def\identityN[#1,#2]{
	_{#1}\mathbb{I}_{#2}^{[N]}
}

\def\identityOne[#1,#2]{
	_{#1}\mathbb{I}_{#2}^{[1^N]}
}

\def\DeltaN[#1]{
	\Delta^{[N]}\left( #1 \right)
}

\def\DeltaOne[#1]{
	\Delta^{[1^N]}\left( #1 \right)
}


%========================================================
% Mathematical Symbols 
\def\half{\frac{1}{2}}
\def\imp{\Longrightarrow}
\def\goto{\rightarrow}
\def\para{\parallel}
\def\vev#1{\langle #1 \rangle}
\def\del{\nabla}
\def\p{\partial}
\newcommand{\bra}[1]{\langle{#1}|}
\newcommand{\ket}[1]{|{#1}\rangle}
\newcommand{\lra}{\leftrightarrow}


% ======================================================
% TEX ENVIRONMENTS
\newenvironment{proposal}[1][Proposal]{\begin{trivlist}
\item[\hskip \labelsep {\bfseries #1}]}{\end{trivlist}}



%=====================================================
% arrows and styles for flavor and gauge nodes
\newcommand{\tikznode}[2]{\relax
\ifmmode%
  \tikz[remember picture,baseline=(#1.base),inner sep=0pt]\node(#1){$#2$};
\else
  \tikz[remember picture,baseline=(#1.base),inner sep=0pt]\node(#1){#2};%
\fi}

\tikzstyle{BFline}=[dashed,black,draw]
\tikzstyle{farrowstyle}=[fcolor,thick,draw]
\tikzstyle{barrowstyle}=[bcolor,thick,draw]

% EXPLICIT QUIVER NODES
\tikzstyle{flavor}=[rectangle,draw=black,thick,inner sep = 0pt, minimum size = 6mm]
\tikzstyle{manifest}=[rectangle,draw=blue!50,thick,inner sep = 0pt, minimum size = 6mm]
\tikzstyle{gauge}=[circle,draw=black,thick,inner sep = 0pt, minimum size = 6mm] %regular gauge node
\tikzstyle{gauge2}=[circle,draw=black!50,thick,inner sep = 0pt, minimum size = 4.5mm] % small circle for SU(N)
\tikzstyle{gauge3}=[rounded rectangle, draw=black!100, thick, minimum size=5mm] % node for long labels

% COMPACT QUIVER NODES
\tikzstyle{gaugefill}=[circle,draw=black, fill = black, thick,inner sep = 0pt, minimum size = 1.5mm]
\tikzstyle{flavorfill}=[rectangle,draw=black, fill = black,,thick,inner sep = 0pt, minimum size = 1.5mm]

% ARROWS
\tikzset{->-/.style={decoration={
  markings,
  mark=at position .5 with {\arrow{>}}},postaction={decorate}}}
\tikzset{-<-/.style={decoration={
  markings,
  mark=at position .5 with {\arrow{<}}},postaction={decorate}}}
\tikzstyle{BFline}=[dashed]
  
% ======================================================
% QUIVER NODES AND ARROWS

% susy nodes and arrows
\def\nodeCS(#1,#2)(#3,#4,#5){ 	% nodeCS(x,y)(nodename,nodelabel, CSlevel)
	\node at (#1,#2) (#3) [gauge,black] {#4};
	\draw[CScolor] (#3) ++(6pt,-7pt) node[anchor=west] {\tiny#5};
}

\def\flavorCS(#1,#2)(#3,#4,#5){ 	% nodeCS(x,y)(nodename,nodelabel, CSlevel)
	\node at (#1,#2) (#3) [flavor,black] {#4};
	\draw[CScolor] (#3) ++(6pt,-8pt) node[anchor=west] {\tiny#5};
}

\def\arrowBF(#1,#2)(#3){ %\arrowBF(node1, node2) (BFlevel)
	\begin{scope}[every node/.style={auto, outer sep=-1pt}]
	\path (#1) edge[->-] node[BFcolor,midway] {\tiny#3} (#2);
	\end{scope}
}

\def\arrowBFlr(#1,#2)(#3,#4){ %\arrowBF(node1, node2) (BFlevel, left/right)
	\begin{scope}[every node/.style={auto=#4, outer sep=-1pt}]
	\path (#1) edge[->-] node[BFcolor,midway] {\tiny#3} (#2);
	\end{scope}
}

\def\dottedBF(#1,#2)(#3){ %\arrowBF(node1, node2) (BFlevel)
	\begin{scope}[every node/.style={auto, outer sep=-1pt}]
	\path (#1) edge[BFline] node[BFcolor,midway] {\tiny#3} (#2);
	\end{scope}
}


% non-susy arrow
\def\farrow(#1,#2,#3){%\farrow(node1, node2, left/right)
	\begin{scope}[every node/.style={auto=#3, outer sep=-1pt}]
	\path (#1) edge[->-,fcolor] node[fcolor,midway] {\tiny$\psi$} (#2);
	\end{scope}
}

\def\barrow(#1,#2,#3){%\farrow(node1, node2, left/right)
	\begin{scope}[every node/.style={auto=#3, outer sep=-1pt}]
	\path (#1) edge[->-,bcolor] node[bcolor,midway] {\tiny$\phi$} (#2);
	\end{scope}
}

\def\barrowBF(#1,#2,#3,#4){%\farrow(node1, node2, left/right)
	\begin{scope}[every node/.style={auto=#3, outer sep=-1pt}]
	\path (#1) edge[->-,bcolor] node[bcolor,midway] {\tiny$\phi$}  node[BFcolor,midway,swap] {\tiny#4} (#2);
	\end{scope}
}

\def\farrowBF(#1,#2,#3,#4){%\farrow(node1, node2, left/right)
	\begin{scope}[every node/.style={auto=#3, outer sep=-1pt}]
	\path (#1) edge[->-,fcolor] node[fcolor,midway] {\tiny$\psi$}  node[BFcolor,midway,swap] {\tiny#4} (#2);
	\end{scope}
}





%===================================================================
% fermionic and bosonic QED
\def\fQED(#1,#2){ %\fQED(Nf,k)
\begin{tikzpicture}{baseline=(current bounding box).center}
	\node at (0,0) (g) [gauge,black] {$1$};
	\draw[CScolor] (g.south east) ++(-3pt,3pt) node[anchor=north west] {\tiny$#2$};
	\node at (1.5,0) (f) [flavor,black] {$#1$};
	\farrow(g,f,left)
\end{tikzpicture}
}

\def\sQED(#1,#2){ %\sQED(Nf,k)
\begin{tikzpicture}{baseline=(current bounding box).center}
	\node at (0,0) (g) [gauge,black] {$1$};
	\draw[CScolor] (g.south east) ++(-3pt,3pt) node[anchor=north west] {\tiny$#2$};
	\node at (1.5,0) (f) [flavor,black] {$#1$};
	\barrow(g,f,left)
\end{tikzpicture}
}
\def\bQED(#1,#2){ \sQED(#1,#2) }







%=====================================================
% TIKZ STYLES
\tikzset{cross/.style={cross out, draw=black, minimum size=5*(#1-\pgflinewidth), inner sep=0pt, outer sep=0pt},
%default radius will be 1pt. 
cross/.default={2pt}}

\tikzset{snake it/.style={decorate, decoration=snake}}

\tikzset{mid arrow/.style={postaction={decorate,decoration={
        markings,
        mark = at position .55 with {\arrow[#1]{Straight Barb[width=5pt]}}
      }}}}
      
\tikzset{mid arrowsm/.style={postaction={decorate,decoration={
        markings,
        mark = at position .55 with {\arrow[#1]{Straight Barb[width=3pt]}}
      }}}}
      
\tikzset{middx arrowsm/.style={postaction={decorate,decoration={
        markings,
        mark = at position .7 with {\arrow[#1]{Straight Barb[width=3pt]}}
      }}}}
      
\tikzset{midsx arrowsm/.style={postaction={decorate,decoration={
        markings,
        mark = at position .4 with {\arrow[#1]{Straight Barb[width=3pt]}}
      }}}}




  
\newcommand{\wigE}{\draw[snake=zigzag]}
\newcommand{\orwigE}{\draw[snake=zigzag,orange]}
%\newcommand{\wigM}{\draw[snake=coil,segment length=3pt]}
\newcommand{\wigM}{\draw[snake=coil,,segment aspect=0,segment length=6pt, double]}
\newcommand{\wigMc}{\draw[->,snake=coil,segment aspect=0,segment length=6pt]}
\newcommand{\wigMB}{\draw[->,snake=zigzag,segment aspect=0,segment length=6pt]}
\newcommand{\wigT}{\draw[dashed]}
\newcommand{\NSbrane}{\draw[ultra thick, blue]}
\newcommand{\Dbrane}{\draw[dashed, ultra thick, gray]}
\newcommand{\Dthree}{\draw[thick, red]}
\newcommand{\hyper}{\draw[-,double distance=3]}
\newcommand{\chir}{\draw[-,mid arrow]}

\newcommand{\convexpath}[2]{
[   
    create hullnodes/.code={
        \global\edef\namelist{#1}
        \foreach [count=\counter] \nodename in \namelist {
            \global\edef\numberofnodes{\counter}
            \node at (\nodename) [draw=none,name=hullnode\counter] {};
        }
        \node at (hullnode\numberofnodes) [name=hullnode0,draw=none] {};
        \pgfmathtruncatemacro\lastnumber{\numberofnodes+1}
        \node at (hullnode1) [name=hullnode\lastnumber,draw=none] {};
    },
    create hullnodes
]
($(hullnode1)!#2!-90:(hullnode0)$)
\foreach [
    evaluate=\currentnode as \previousnode using \currentnode-1,
    evaluate=\currentnode as \nextnode using \currentnode+1
    ] \currentnode in {1,...,\numberofnodes} {
  let
    \p1 = ($(hullnode\currentnode)!#2!-90:(hullnode\previousnode)$),
    \p2 = ($(hullnode\currentnode)!#2!90:(hullnode\nextnode)$),
    \p3 = ($(\p1) - (hullnode\currentnode)$),
    \n1 = {atan2(\y3,\x3)},
    \p4 = ($(\p2) - (hullnode\currentnode)$),
    \n2 = {atan2(\y4,\x4)},
    \n{delta} = {-Mod(\n1-\n2,360)}
  in 
    {-- (\p1) arc[start angle=\n1, delta angle=\n{delta}, radius=#2] -- (\p2)}
}
-- cycle
}



% COMPACT FUNDAMENTAL BLOCKS
\def\simpFund[#1]{ % #1 is N
	\begin{tikzpicture}
        \node at (0,0) (g) [flavor, black] {$#1$};
        \node at (0,1.5) (f) [flavor,black] {$1$};
        \draw[CScolor] (g)+(0,-.6) node {$_{( \text{-}\frac{1}{2}, \text{-}\frac{1}{2})}$};
        %\draw[CScolor,right] (f2)+(.2,-.4) node {$_{\text{-}\frac{#1}{2}}$};
        %\draw[BFcolor,left] (-.1,.75) node {$_{+1}$};
        \draw[->-] (g)--(f);
        \draw[right] (g)+(.4,0) node {$\mathbb{I}^{[#1]}$};
	\end{tikzpicture}
}

\def\simpAFund[#1]{ % #1 is N
	\begin{tikzpicture}
        \node at (0,0) (g) [flavor, black] {$#1$};
        \node at (0,1.5) (f) [flavor,black] {$1$};
        \draw[CScolor] (g)+(0,-.6) node {$_{(\text{-}\frac{1}{2}, \text{-}\frac{1}{2})}$};
        %\draw[CScolor,right] (f2)+(.2,-.4) node {$_{\text{-}\frac{#1}{2}}$};
        %\draw[BFcolor,left] (-.1,.75) node {$_{+1}$};
        \draw[->-] (f)--(g);
        \draw[right] (g)+(.4,0) node {$\mathbb{I}^{[#1]}$};
	\end{tikzpicture}
}

\def\simpFundNMirror[#1]{ % #1 is N
	\begin{tikzpicture}
		\node at (-1.5,0) (f1) [flavor,black] {$#1$};
        \node at (0,0) (g1) [gauge, black] {$1^#1$};
        \node at (1,1) (g2) [gauge,black] {$1^#1$};
        \node at (2.5,1) (f2) [flavor,black] {$#1$};
        \draw[->-] (g1)--(g2);
        \draw[decorate,decoration={coil,segment length=4pt}] (f1)--(g1) node[midway,above] {$+$};
        \draw[decorate,decoration={coil,segment length=4pt}] (g2)--(f2)node[midway,above] {$-$};
	\end{tikzpicture}
}

\def\simpAFundNMirror[#1]{
	\begin{tikzpicture}
		\node at (-1.5,1) (f1) [flavor,black] {$#1$};
        \node at (0,1) (g1) [gauge, black] {$1^#1$};
        \node at (1,0) (g2) [gauge,black] {$1^#1$};
        \node at (2.5,0) (f2) [flavor,black] {$#1$};
        \draw[->-] (g2)--(g1);
        \draw[decorate,decoration={coil,segment length=4pt}] (f1)--(g1) node[midway,above] {$+$};
        \draw[decorate,decoration={coil,segment length=4pt}] (g2)--(f2)node[midway,above] {$-$};
	\end{tikzpicture}
}

\def\simpBifLNN[#1]{
	\begin{tikzpicture}
        \node at (0,0) (g) [flavor, black] {$#1$};
        \node at (1.5,0) (f) [flavor,black] {$#1$};
        \draw[CScolor] (g)+(0,-.6) node {$_{(\text{-}\frac{#1}{2},\text{-}\frac{#1}{2})}$};
        \draw[CScolor] (f)+(0,-.6) node {$_{(\text{-}\frac{#1}{2},\text{-}\frac{#1}{2})}$};
        \draw[BFcolor] (.75,.3) node {$_{+1}$};
        \draw[->-] (f)--(g);
	\end{tikzpicture}
}

\def\simpBifRNN[#1]{
	\begin{tikzpicture}
		\node at (0,0) (g) [flavor, black] {$#1$};
        \node at (1.5,0) (f) [flavor,black] {$#1$};
        \draw[CScolor] (g)+(0,-.6) node {$_{(\text{-}\frac{#1}{2},\text{-}\frac{#1}{2})}$};
        \draw[CScolor] (f)+(0,-.6) node {$_{(\text{-}\frac{#1}{2},\text{-}\frac{#1}{2})}$};
        \draw[BFcolor] (.75,.3) node {$_{+1}$};
		\draw[->-] (g)--(f);
	\end{tikzpicture}
}

\def\simpBifLNNMirror[#1]{
\;
	\begin{tikzpicture}
		\node at (-1.5,0) (f1) [flavor,black] {$#1$};
        \node at (0,0) (g) [gauge, black] {$1^#1$};
        \node at (1.5,0) (f2) [flavor,black] {$#1$};
        \node at (1,-1) (f3) [flavor,black] {$1$};
        \draw[decorate,decoration={coil,segment length=4pt}] (f1)--(g) node[midway,above] {$+$};
        \draw[decorate,decoration={coil,segment length=4pt}] (g)--(f2)node[midway,above] {$-$};
		\draw[->-] (f3)--(g);
	\end{tikzpicture}
\;
}

\def\simpBifRNNMirror[#1]{
\;
	\begin{tikzpicture}
		\node at (-1.5,0) (f1) [flavor,black] {$#1$};
        \node at (0,0) (g) [gauge, black] {$1^#1$};
        \node at (1.5,0) (f2) [flavor,black] {$#1$};
        \node at (-1,1) (f3) [flavor,black] {$1$};
        \draw[decorate,decoration={coil,segment length=4pt}] (f1)--(g) node[midway,above] {$+$};
        \draw[decorate,decoration={coil,segment length=4pt}] (g)--(f2)node[midway,above] {$-$};
		\draw[->-] (g)--(f3);
	\end{tikzpicture}
\;
}

\def\simpGenBifund[#1]{ % #1 is N
\;
	\begin{tikzpicture}[baseline=0]
		    \node at (0,0) (g1) [flavor, black] {$#1$};
		    \node at (1.3,0) (f1) [flavor,black] {$1^#1$};
		    \draw[decorate,decoration={coil,segment length=4pt}] (f1)--(g1);
	\end{tikzpicture}
}

\def\simpGenBifundA[#1]{ % #1 is N
\;
	\begin{tikzpicture}[baseline=0]
		    \node at (0,0) (g1) [flavor, black] {$1^#1$};
		    \node at (1.3,0) (f1) [flavor,black] {$#1$};
		    \draw[decorate,decoration={coil,segment length=4pt}] (f1)--(g1);
	\end{tikzpicture}
}

\def\simpBifLOne[#1]{
\;
	\begin{tikzpicture}[baseline=10]
		    \node at (1,0) (g) [flavor, black] {$#1$};
		    \node at (0,1) (f) [flavor,black] {$#1$};
		    \draw[->-] (g)--(f);
	\end{tikzpicture}
\;
}

\def\simpBifROne[#1]{
\;
	\begin{tikzpicture}[baseline=10]
		    \node at (0,0) (g) [flavor, black] {$#1$};
		    \node at (1,1) (f) [flavor,black] {$#1$};
		    \draw[->-] (g)--(f);
	\end{tikzpicture}
\;
}

\def\simpFundOneMirror[#1]{
\;
	\begin{tikzpicture}[baseline=0]
		\node at (-1.5,0) (f1) [flavor,black] {$1^#1$};
		    \node at (0,0) (g1) [gauge, black] {$#1$};
      \draw[red] (-.5,-.6) node {$_{(-\frac{1}{2}, #1-\frac{1}{2})}$};
        \draw[BFcolor] (.5,.5) node {$_1$};
      \draw[red] (1.5,-.6) node {$_{(-\frac{1}{2}, #1-\frac{1}{2})}$};
		    \node at (1,0) (g2) [gauge,black] {$#1$};
		    \node at (2.5,0) (f2) [flavor,black] {$1^#1$};
		    \draw[->-] (g1)--(g2);
		    \draw[decorate,decoration={coil,segment length=4pt}] (f1)--(g1) node[midway,above] {$+$};
			\draw[decorate,decoration={coil,segment length=4pt}] (g2)--(f2)node[midway,above] {$-$};
	\end{tikzpicture}
\;
}

\def\simpFundOne[#1]{
\;
	\begin{tikzpicture}[baseline=0]
		    \node at (0,0) (g1) [flavor, black] {$1^#1$};
		    \node at (0,1.5) (f2) [flavor,black] {$1$};
		    \draw[->-] (g1)--(f2);
	\end{tikzpicture}
\;
\mathbb{I}^{[1^#1]}
\;
}

\def\simpAFundOneMirror[#1]{
\;
	\begin{tikzpicture}[baseline=0]
		\node at (-1.5,0) (f1) [flavor,black] {$1^#1$};
		    \node at (0,0) (g1) [gauge, black] {$#1$};
      \draw[red] (-.5,-.6) node {$_{(-\frac{1}{2}, #1-\frac{1}{2})}$};
        \draw[BFcolor] (.5,.5) node {$_1$};
      \draw[red] (1.5,-.6) node {$_{(-\frac{1}{2}, #1-\frac{1}{2})}$};
		    \node at (1,0) (g2) [gauge,black] {$#1$};
		    \node at (2.5,0) (f2) [flavor,black] {$1^#1$};
		    \draw[->-] (g2)--(g1);
		    \draw[decorate,decoration={coil,segment length=4pt}] (f1)--(g1) node[midway,above] {$+$};
			\draw[decorate,decoration={coil,segment length=4pt}] (g2)--(f2)node[midway,above] {$-$};
	\end{tikzpicture}
\;
}

\def\simpAFundOne[#1]{
\;
	\begin{tikzpicture}[baseline=0]
		    \node at (0,0) (g1) [flavor, black] {$1^#1$};
		    \node at (0,-1.5) (f2) [flavor,black] {$1$};
		    \draw[->-] (f2)--(g1);
	\end{tikzpicture}
\;
\mathbb{I}^{[1^#1]}
\;
}

\def\simpBifLOneMirror[#1]{
\;
	\begin{tikzpicture}[baseline=0]
		\node at (-1.5,0) (f1) [flavor,black] {$1^#1$};
		    \node at (0,0) (g) [gauge, black] {$#1$};
                \draw[red] (0,-.6) node {$_{(-\frac{1}{2}, #1-\frac{1}{2})}$};
		    \node at (1.5,0) (f2) [flavor,black] {$1^#1$};
		 	 \node at (-1.5,1) (f3) [flavor,black] {$1$};
		    \draw[decorate,decoration={coil,segment length=4pt}] (f1)--(g) node[midway,above] {$+$};
			\draw[decorate,decoration={coil,segment length=4pt}] (g)--(f2)node[midway,above] {$-$};
		\draw[->-] (f3)--(g);
	\end{tikzpicture}
\;
}

\def\simpBifROneMirror[#1]{
\;
	\begin{tikzpicture}[baseline=0]
		\node at (-1.5,0) (f1) [flavor,black] {$1^#1$};
		    \node at (0,0) (g) [gauge, black] {$#1$};
         \draw[red] (-.5,-.6) node {$_{(-\frac{1}{2}, #1-\frac{1}{2})}$};
		    \node at (1.5,0) (f2) [flavor,black] {$1^#1$};
		 	 \node at (1.5,-1) (f3) [flavor,black] {$1$};
		    \draw[decorate,decoration={coil,segment length=4pt}] (f1)--(g) node[midway,above] {$+$};
			\draw[decorate,decoration={coil,segment length=4pt}] (g)--(f2)node[midway,above] {$-$};
		\draw[->-] (g)--(f3);
	\end{tikzpicture}
\;
}
%%%%%%%%%%%%%%%%%%%%%%%%%%%%%%%%%%%%%%%%%%%%%%%%%%%%%%%%%%%%%%%

% Zones
% 1
% 1t
% 2
 % 2t 
 % 3
 % 3t
 % 22t - 2 and 2t
 % 33t - 3 and 3t
 % N0
 % mchi
 % k0
 % F0
% KF planes
\def\KFplane(#1)(#2,#3){ % #1=zone, #2=scale #3
%\begin{tikzpicture}[scale=#2]
\begin{scope}[scale=#2]
% colored zones
% Zone 1
\ifthenelse{\equal{#1}{1}}{
    \fill[fill=#3] (4,0)--(6,0)--(6,2)--cycle;
}{
% Zone 1t
\ifthenelse{\equal{#1}{1t}}{
    \fill[fill=#3] (4,0)--(2,0)--(2,2)--cycle;
}{
% Zone 2
\ifthenelse{\equal{#1}{2}}{
    \fill[fill=#3] (4,0)--(6,2)--(6,6)--(4,4)--cycle;
}{
% Zone 2t
\ifthenelse{\equal{#1}{2t}}{
    \fill[fill=#3] (4,0)--(4,4)--(2,2)--cycle;
}{
% Zone 3
\ifthenelse{\equal{#1}{3}}{
    \fill[fill=#3] (4,4)--(6,6)--(2,6)--cycle;
}{
% Zone 3t
\ifthenelse{\equal{#1}{3t}}{
    \fill[fill=#3] (0,4)--(2,2)--(4,4)--(2,6)--(0,6)--cycle;
}{
% Zone 2 and 2t
\ifthenelse{\equal{#1}{22t}}{
    \fill[fill=#3] (4,0)--(6,2)--(6,6)--(4,4)--cycle;
    \fill[fill=#3] (4,0)--(4,4)--(2,2)--cycle;
}{
% Zone 3 and 3t
\ifthenelse{\equal{#1}{33t}}{
    \fill[fill=#3] (4,4)--(6,6)--(2,6)--cycle;
     \fill[fill=#3] (0,4)--(2,2)--(4,4)--(2,6)--(0,6)--cycle;
}
}}}}}}}


% Axes
\draw[->,thick] (0,0) -- (6,0) node[right] {$\scriptstyle F$};
\draw[->,thick] (0,0) -- (0,6) node[above] {$\scriptstyle -2k$};

% Labels on axes
\node[left] at (0,4) {$\scriptstyle 2N$};
\node[below] at (2,0) {$\scriptstyle N$};
\node[below] at (4,0) {$\scriptstyle 2N$};

% Black zone boundaries
\draw[thick] (2,0) -- (2,2);
\draw[thick] (4,0) -- (0,4);
\draw[thick] (4,0) -- (4,4);
\draw[thick] (4,4)--(2,6);
\draw[thick] (2,2)--(6,6);
\draw[thick] (4,0)--(6,2);

%highlighted lines
% Zone 1
\ifthenelse{\equal{#1}{N4}}{
    \draw[#3,line width=3pt] (4,0)--(6,2);
}{
% Zone 1t
\ifthenelse{\equal{#1}{k0}}{
    \draw[#3,line width=3pt] (4,0)--(6,0);
}{
% Zone 2
\ifthenelse{\equal{#1}{mchi}}{
    \draw[#3,line width=3pt] (2,2) -- (6,6);
}{
% Zone 2t
\ifthenelse{\equal{#1}{F0}}{
    \draw[#3,line width=3pt] (0,4)--(0,6);
}}}}
\end{scope}
%\end{tikzpicture}
}

  






